\section{Fear of preventable harm: the errors we can no longer accept}
\label{sec:fear}

The second appeal is the fear of harm that did not have to happen. A generation of
health policy, from drug-safety regulation to device oversight to electronic
health record requirements, was built on stories of missed diagnoses, medication
errors, unsafe products, and delayed treatment. The fear is not abstract. The
landmark report \textit{To Err Is Human} estimated that tens of thousands of
Americans die each year from preventable medical error \cite{kohn2000toerr}, and a
later analysis ranked medical error among the leading causes of death in the
United States \cite{makary2016medical}. The underlying message of this pillar is
plain: without a better check, more people will be harmed than need to be.

Why does error accumulate? Because the legacy clinical-trial and peer-review
system passes work through many manual handoffs, and each handoff is an
independent opportunity for a mistake that compounds downstream. The verified path
inserts an automated gate at each step, so a defect is caught where it occurs
rather than discovered after a patient is affected. Figure 4 draws the
contrast, and the engineering record behind it is concrete: a surgical-humanoid
assurance run passed 172 of 172 automated tests, 64 of them inside a ten-gate
suite, before any action was permitted \cite{Kawchak2026MobilePancreaticCancerUnitree}.

\begin{narrativefig}
\node[nfrisk] (a) at (0,0) {Manual handoff};
\node[nfrisk] (b) at (4,0) {Error compounds};
\node[nfrisk] (c) at (8,0) {Harm reaches patient};
\node[nftwo] (d) at (1.7,-1.6) {Generate};
\node[nfdec] (e) at (5.2,-1.6) {Gate};
\node[nfhope] (f) at (8.6,-1.6) {Defect contained};
\draw[nfedge] (a)--(b); \draw[nfedge] (b)--(c);
\draw[nfedge] (d)--(e); \draw[nfedge] (e)--(f);
\end{narrativefig}
\vspace{-0.7cm}
\figcaption{Figure 4. Compounding error vs.\ a contained defect, below. The gate stops a defect at its source.}

\autoref{tab:fear} sets the two paths side by side. The comparison is the kind of
scenario-based risk information that helps a committee visualize a future loss and
turn a what-if into a clear policy goal.

\needspace{9\baselineskip}\begin{table}[H]
\footnotesize
\begin{tabularx}{\textwidth}{@{}L{4.95cm} Y Y@{}}
\toprule
\textbf{Step} & \textbf{Legacy path} & \textbf{Verified path}\\
\midrule
Where a defect is caught & After the fact, if at all & At the gate where it occurs\\
Who must notice it & A tired human, unaided & An automated check, every time\\
Record of the decision & Often informal & Hash-chained and auditable\\
Effect of one error & Compounds downstream & Contained and escalated\\
\bottomrule
\end{tabularx}
\caption{The legacy path and the verified path, step by step.}
\label{tab:fear}
\end{table}

A bill cannot promise that nothing will ever go wrong. It can promise that the
system checks itself at every step, so the same mistake is not made twice and is
never made invisibly. When committee \emph{markup} sharpens this language, the
safety guarantee becomes part of the permanent record. The quiet question returns:
why did we ever fully trust a process with so many chances for compounding human
error?
